\chapter{INTRODUCTION}

\section{Background of AI in Recruitment}
Technology has drastically changed the HR sector, among many sectors. One way technology has changed HR is the development of websites where job seekers find jobs. Instead of job hunting locally as they have for many years, they can apply to virtually any posting around the globe. With this newfound ease of access, job seekers apply to more postings, and HR departments must now deal with more job postings and a higher volume of applicants. In fact, most tech companies report receiving thousands of postings just a couple of hours after posting a job. Because of this, many companies resort to a highly automated hiring process. This system also uses a new generation of Applicant Tracking Systems (ATS). This modern generation of ATS has upgraded from a primitive system of storing CVs to an automated hiring process, where ATS will scan, assess, and eliminate applicants who would likely not increase the productivity of the company.

\section{Problems in Manual Resume Screening}
There exist three main issues associated with manual resume screening: the use of too much time, exhaustion of resources, and bias. Often, HR managers have to scan huge amounts of applicant data in only a couple of seconds per individual resume. Within such a limited period of time, one can identify only certain visual elements, the names of certain universities or previous employers of the applicants, thus reducing the quality of assessment. It cannot be ruled out that some kind of a bias will occur during rapid mental analysis, whether demographic or 'layout bias'/'format bias'. Under layout bias conditions, the candidate is appreciated or discredited solely because of the visual appearance of the document or complexity of its design. Finally, the economic cost of employing huge recruitment teams that perform only the first stage of assessment is totally unacceptable for SMEs.

\section{The Limitations of Keyword Matching}
In the first generation of these ATS software systems, everything was based purely on rigid Boolean keyword searching. When the JD asked for "Python" and "Data Analysis", the ATS software blindly searched through the applicant's uploaded resume for the presence of those exact strings of characters. There were numerous drawbacks associated with this strategy that were directly detrimental to the recruitment process. First, this system fails to account for the "Implicit Competency" problem - someone writing "I architected scalable backend web systems" has demonstrated excellent competency as an engineer, but if the ATS is specifically looking for the string of keywords "backend developer," then the applicant is out of luck. Finally, this early form of keyword filtering is highly vulnerable to "keyword stuffing". Keyword stuffing is the technique of copying the whole JD description into the white-colored margins of one's resume.

\section{The Need for Explainable and Cost-Efficient Screening}
However, once it became clear to the industry that there were inherent limitations to the old keyword system, the industry quickly shifted toward very complicated deep learning models, which employed BERT (Bidirectional Encoder Representations from Transformers) and other LLMs (Large Language Models) to do semantic matching. Although this approach solves the problem of keyword mismatch by recognizing that ``REST'' and ``RESTful'' are semantically similar, it creates two major new issues: sky-high computing costs and a total lack of transparency. Most enterprise AI tools operate as ``black boxes.'' In this case, even when a deep neural net ranks the candidate at 92\% match, the recruiter does not have any mathematical or logical explanation for why it is so. And in the event that the rejected candidate asks for an explanation or takes the employer to court, it will be impossible to audit the black box of the neural net to prove that the candidate was truly selected based on merit only. Therefore, there is a great demand on the market for AI recruiting solutions that can provide full auditability and transparency, as well as low computational costs.

\section{Motivation for the Project}
A primary motivation for this project can be found in the fact that despite increasing automation of the hiring process with advanced technologies, most options offered to SMEs are either inadequate legacy solutions (ATS) or overly costly and difficult to understand (LLMs). The current state of affairs shows that the process of hiring employees continues to consume resources, and manual review of resumes leads to inconsistent selection of candidates. Hence, there is a definite need to design an approach to improve on the idea of keyword matching using intelligent parsing and math-based bounds without resorting to extensive calculations associated with generative AI.

\section{System Boundaries and Functional Constraints}
It is crucial to set the proper functional scope of the developed solution to avoid overpromising. The system under development cannot be regarded as an ideal expensive commercial AI service with profound capabilities of semantic analysis. It certainly is not the ultimate answer to one of the most complicated sociological problems related to recruiting people. Finally, it should not replace human interviews that are a must in any recruiting practice. What this project offers is a fast lightweight, easy-to-understand solution designed to do simple pre-selection. This is done at the expense of depth of semantic analysis in favor of speed, complete explainability, and zero expenses on server hosting.

\section{Problem Statement}
Besides being a highly expensive and time-consuming procedure due to the significant volume of applications submitted by the candidates for any single job opening, the problem with current practice is that the recruitment process often encounters inconsistency and bias issues linked to manual resume screening. At present, even state-of-the-art technologies offered by modern AI for businesses are still limited in their capabilities, conducting rather primitive analysis of the candidates' documents using keyword search methods, while more advanced semantic AI solutions fail due to the lack of decision explanation possibilities. As a result, there arises a need for an automated candidate-to-position matching method that is cost-efficient and speedy as well as fully transparent for users.

\section{Organization of the Report}
This paper is organized as follows. The second chapter gives a brief overview of the relevant literature. Project goals and objectives are discussed along with the research gaps filled in Chapter 3. Chapter 4 is focused on the description of system design and the algorithms used. Software implementation and technical stack are provided in Chapter 5. Experiments conducted during prototype evaluation are reported in Chapter 6. The results obtained during benchmarking are presented in Chapter 7. Limitations of the developed prototype are discussed in Chapter 8.
